\begin{frame}{학습된 벡터의 놀라운 성질}
  \begin{center}
    \[
      \text{vec}(\text{king}) - \text{vec}(\text{man})
      + \text{vec}(\text{woman})
      \;\approx\;
      \text{vec}(\text{queen})
    \]
  \end{center}
  \begin{itemize}
    \item ``왕에서 남자다움을 빼고 여자다움을 더하면 여왕'' --- 학습
      당시 \textbf{전혀 명시적으로 가르치지 않았는데도} 성립.
    \item 벡터 공간의 \textbf{기하학적 구조}가 스스로 그렇게 정리된
      것.
  \end{itemize}
  \vskip0.6em
  {\footnotesize (Mikolov, Sutskever, Chen, Corrado \& Dean, 2013 ---
  단어 유사도·유추(analogy) 분석을 본격 도입한 논문.)}
\end{frame}
